\section{Localized energy and individual coefficient bounds}\label{sec:energy}

The first estimate uses the full spatial gradient only where the
eigenfunction equation fails. After smoothing and cutting off near
the vertices, this will confine the resulting error to small cone
annuli. The zero-frequency version of this argument appears in
\cite[Sections~3--4]{Companion}; the spectral correction below is
needed when \(\lambda\ne0\). The equal complex-derivative norms and
coefficient recurrences also occur in Forni and Moll's Sobolev
arguments \cite[Lemma~3.3, Corollary~4.2, and Lemma~5.1]{FM}.
Twisted complex-derivative energy identities provide a related
spectral precedent in Forni's work \cite[Section~3]{Forni2022}.

\begin{lemma}[Localized coefficient estimate]\label{lem:energy}
Let \(w\in L^2(SM)\) have compact spatial support in \(M\) and be
smooth in position with values in angular \(L^2\) in every flat chart.
Suppose \(Aw\) vanishes over \(M\setminus E\), where \(E\subset M\)
is measurable. Define
\(D_j=\norm{\pa w_j}_2^2\). Then, for every \(m\in\Z\),
\begin{equation}\label{eq:localized-energy}
 \tfrac12D_m+D_{m+1}
 \le \norm{Aw}_2\norm{\ind_E\nabla w}_2
       +\tfrac{\lambda^2}{2}\norm{w_{m+1}}_2^2,
 \qquad
 D_j=\tfrac14\norm{\nabla w_j}_2^2.
\end{equation}
\end{lemma}

\begin{proof}
Compact support in the regular surface permits integration by parts
between the coefficient bundles. The local formal adjoints are
\(\pa^*=-\pb\) and \(\pb^*=-\pa\); the derivatives commute in
flat charts. Hence
\begin{equation}\label{eq:complex-energies}
 \norm{\pa w_j}_2^2=\norm{\pb w_j}_2^2
                 =\tfrac14\norm{\nabla w_j}_2^2.
\end{equation}
These identities are intrinsic, since all transition factors are
constant and unitary. In particular, no integration takes place across
a cone point.

Put
\(C_j=\operatorname{Re}\ip{i\lambda w_{j+1}}{\pa w_j}\).
Equation~\eqref{eq:modes} gives
\begin{equation}\label{eq:spectral-telescope}
 D_j-D_{j+2}
 =\operatorname{Re}\ip{(Aw)_{j+1}}{\pa w_j-\pb w_{j+2}}
   +C_j-C_{j+1}.
\end{equation}
For the spectral term, the sign follows from
\[
 \operatorname{Re}\ip{i\lambda w_{j+1}}{\pb w_{j+2}}
 =\operatorname{Re}\ip{-i\lambda\pa w_{j+1}}{w_{j+2}}
 =C_{j+1}.
\]
The correction \(C_j-C_{j+1}\) couples consecutive indices, so a sum
over one parity class would not telescope. Sum
\eqref{eq:spectral-telescope} over \emph{every} integer
\(j=m,\ldots,N\). The left side is
\(D_m+D_{m+1}-D_{N+1}-D_{N+2}\), and the last two terms on the
right telescope to \(C_m-C_{N+1}\). All high endpoints tend to zero:
the coefficients of \(w\) and their spatial derivatives are square
summable, by the hypotheses on \(w\). The remaining sum of pairings
converges absolutely by Cauchy--Schwarz. Since \(Aw\) vanishes off
the part over \(E\), its absolute value is at most
\[
 \norm{Aw}_2
 \left(\int_E\sum_{j\ge m}
       |\pa w_j-\pb w_{j+2}|^2\,d\bar A\right)^{1/2}
 \le \norm{Aw}_2\norm{\ind_E\nabla w}_2.
\]
The last inequality uses the pointwise bound
\[
 \sum_{j\ge m}|\pa w_j-\pb w_{j+2}|^2
 \le 2\sum_{j\in\Z}(|\pa w_j|^2+|\pb w_j|^2)
 =\sum_{j\in\Z}|\nabla w_j|^2.
\]
Finally,
\[
 |C_m|\le |\lambda|\norm{w_{m+1}}_2\sqrt{D_m}
 \le \tfrac12D_m+\tfrac{\lambda^2}{2}\norm{w_{m+1}}_2^2,
\]
which proves \eqref{eq:localized-energy}.
\end{proof}

We now construct regularizations to which Lemma~\ref{lem:energy}
applies. Write \(r(x)=d(x,V)\) in \(\overline M\), also using \(r\)
for its pullback to \(SM\). Near each cone point, an embedded flat
disk about \(x\) has radius at least \(c r(x)\), for a fixed
\(c>0\): this follows by choosing uniformly sized disks on the compact
circle at radius one in the model cone and scaling. Compactness away
from the tips therefore gives a constant \(K>2\) such that, for all
sufficiently small \(\eps>0\), every point with
\(r>(K-2)\eps\) has an embedded flat disk of radius \(4\eps\).
For each fixed \(L>0\), the cone area formula also gives
\begin{equation}\label{eq:cone-area}
 \bar A\{r\le L\eps\}=O_L(\eps^2).
\end{equation}

Choose a nonnegative radial function
\(\rho\in C_c^\infty(\R^2)\), supported in the unit disk and with
Euclidean integral one, and set
\(\rho_\eps(y)=\eps^{-2}\rho(y/\eps)\). For \(r(x)>K\eps\),
define spatial convolution at fixed parallel direction by
\begin{equation}\label{eq:smoothing-definition}
 (S_\eps u)(x,\theta)
 =\int_{\R^2}\rho_\eps(y)u(x-y,\theta)\,dy
\end{equation}
in the embedded flat disk. Radiality makes this definition independent
of its oriented chart. Unless a derivative is being taken, extend
\(S_\eps u\) by zero to \(r\le K\eps\).

\begin{lemma}[Spatial regularization]\label{lem:smoothing}
For every integer \(\ell\ge0\), the preceding operators satisfy
\begin{equation}\label{eq:smoothing-bounds}
 \eps^\ell\norm{\nabla^\ell S_\eps u}_{L^2(r>K\eps)}
 \le C_\ell\norm{u}_2.
\end{equation}
Their derivatives here are taken before zero extension. The output
is smooth in position with values in angular \(L^2\), and
\(S_\eps u\to u\) in \(L^2(SM)\) for every fixed \(u\in L^2(SM)\).
For bounded input,
\[
 \norm{S_\eps u}_\infty\le\norm{u}_\infty,
 \qquad
 \eps^\ell\norm{\nabla^\ell S_\eps u}_{L^\infty(r>K\eps)}
 \le C_\ell\norm{u}_\infty.
\]
Convolution commutes locally with \(X,Y\), and \(A\), in
distributions. If \(L>K+4\) and \(g=0\) on \(r<L\eps\), then
\(S_\eps g=0\) on \(r<(L-1)\eps\); its zero extension has compact
spatial support in \(M\) and creates no boundary derivative.
In particular, when also \(Ag\in L^2\),
\begin{equation}\label{eq:smoothing-commutes}
 A(S_\eps g)=S_\eps(Ag)
\end{equation}
globally on the regular bundle, with the same zero-extension convention.
\end{lemma}

\begin{proof}
In a common flat neighborhood of nearby centers, this is ordinary
convolution with a kernel whose derivatives have the usual scaling.
Cauchy--Schwarz bounds the squared angular \(L^2\) norm of
\(\eps^\ell\nabla^\ell S_\eps u\) at a center by
\(C_\ell\eps^{-2}\) times the integral of
\(\norm{u(y,\cdot)}_{L^2_\theta}^2\) over its \(\eps\)-disk.
Parallel identification of tangent circles preserves angular measure.
A point in such a disk satisfies \(r>(K-1)\eps\), and all centers
which can contribute to it lie in a disk of area \(O(\eps^2)\).
Integrating the bound over centers proves
\eqref{eq:smoothing-bounds}. The same kernel estimates give the
\(L^\infty\) bounds; nonnegativity and mass one give the contraction
bound at order zero. Constant-coefficient differentiation in each
chart proves local commutation and spatial smoothness.

For a smooth compactly supported input on \(SM\), the usual
approximate-identity argument gives \(L^2\) convergence, including
the zero extension for sufficiently small \(\eps\). Such inputs
are dense in \(L^2(SM)\), so the uniform bound at order zero proves
convergence for every fixed input.

The function \(r\) is 1-Lipschitz. If \(r(x)<(L-1)\eps\), every
point sampled by its convolution disk has \(r<L\eps\), proving
the support claim. The convolution consequently vanishes in a whole
neighborhood of the boundary \(r=K\eps\), so zero extension produces
no boundary distribution. Local commutation gives
\eqref{eq:smoothing-commutes} there; on the omitted region both sides
vanish, since \(Ag\) has no support where \(g\) vanishes.
\end{proof}

\begin{lemma}[Individual coefficient regularity]\label{lem:bounded-modes}
Suppose \(h\in L^\infty(SM)\), \(\norm{h}_\infty\le H\), and
\((X-i\lambda)h=0\) locally in distributions on the regular bundle.
Every Fourier coefficient then has a global square-integrable flat
gradient, and
\begin{equation}\label{eq:bounded-modes}
 \norm{\nabla h_j}_2\le C_\lambda H\qquad(j\in\Z).
\end{equation}
The constant depends only on the fixed triangle and \(\lambda\).
\end{lemma}

\begin{proof}
Choose a smooth spatial cutoff \(\chi_\eps\), equal to zero on
\(r\le(K+1)\eps\), equal to one on \(r\ge(K+2)\eps\), and
satisfying \(0\le\chi_\eps\le1\) and
\(|\nabla\chi_\eps|\le C/\eps\). Its transition region
\(E_\eps\) consists of three disjoint cone annuli for small
\(\eps\). The field
\(w_\eps=\chi_\eps S_\eps h\), extended by zero near the tips,
has compact spatial support and the smoothness required by
Lemma~\ref{lem:energy}. Local commutation and \(Ah=0\) give
\[
 Aw_\eps=(X\chi_\eps)S_\eps h.
\]
This field is supported over \(E_\eps\), whose area is
\(O(\eps^2)\). On those annuli Lemma~\ref{lem:smoothing} gives
\[
 |Aw_\eps|+|\nabla w_\eps|\le CH/\eps.
\]
Consequently,
\(\norm{Aw_\eps}_2\le CH\) and
\(\norm{\ind_{E_\eps}\nabla w_\eps}_2\le CH\).
Also \(\norm{(w_\eps)_{j+1}}_2\le\norm{w_\eps}_2\le H\).
Equation~\eqref{eq:localized-energy} therefore bounds
\(\norm{\nabla(w_\eps)_j}_2\) by \(C_\lambda H\), independently
of \(\eps\) and \(j\).

We have \(w_\eps\to h\) in \(L^2\): use
Lemma~\ref{lem:smoothing} and the shrinking area of the region where
\(\chi_\eps\ne1\). For each fixed \(j\), choose a weakly convergent
subsequence of the globally bounded coefficient gradients. Testing in
regular charts identifies its limit as the distributional gradient
of \(h_j\). Weak lower semicontinuity proves
\eqref{eq:bounded-modes}.
\end{proof}

This estimate controls one coefficient at a time; it does not yet
give summability of their energies over \(j\). Its hypothesis is only
the local equation \(Ah=0\), so it will apply to weak limits without
requiring a separate flow-invariance argument. The next section uses
long orbit averages to obtain the additional identities needed for
the sum over all coefficients.
